\section{Fabre d'Olivet on the Borean Race}

An early review of two books by Fabre d'Olivet, translated into English, was published in New York Tribune, 7 Aug 1921\footnote{\url{http://chroniclingamerica.loc.gov/lccn/sn83030214/1921-08-07/ed-1/seq-50/;words=Fabre+Borean+D\%27Olivet+d+Olivet}}.

Besides his monumental hermeneutical study of the history of the Borean race, Fabre d'Olivet wrote an interpretation of the Golden Verses of \textbf{Pythagoras}. Born some 800 years after \textbf{Akhnaton}, Pythagoras was an initiate in Egypt and brought monotheism into Greece. While outwardly participating in the polytheistic rites of the city, the initiates of the Pythagorean school were secretly following a different path.

Monotheism and the Logos were taught in the ancient mystery schools. Later, Socrates was a bit too careless and was condemned to death in Athens for the crime of impiety. Nevertheless, these teachings were developed by Plato and Aristotle. However, it was the arrival of Christianity that made the esoteric teaching of monotheism as part of its exoteric teaching … for better or worse. However, as we shall see in future posts, it was a necessary development in the Destiny of the Borean race, as it was more suited to Empire than the ancestor worship of the city-states.

The text of the review follows:

\paragraph{Pythagoras Revealed Through Medium of a French Disciple}
\subparagraph{New Light Cast on Theories of Greek Sage Who Believed in Supreme Deity 500 Years Before Christ}
By Grace Phelps

\emph{THE GOLDEN VERSES OP PYTHAGORAS} By Fabre d'Olivet. Done into English by Nayan Louise Redfield.

Published by O. P. Putnam'a Sons.

\emph{HERMENEUTIC INTERPRETATION OF THE ORIGIN OF THE SOCIAL STATE OF MAN}. By Fabre d Olivet.

Done into English by Nayan Louise Redfield. Published by G. P. Putnam's Sons.

``What's the World Coming To?" is a question that the moderns are not alone in asking. The ancients discussed the subject rather fully, and when one of their wise men thought he had the answer he formed a cult and instructed his followers in a philosophy that enabled them to bear up under the social inequalities and various forms of injustice of the day. Today we translate the ancients.

Pythagoras was one such sage and D'Olivet, his French translator, was another, but their teachings did not endear them to the people of their time. Pythagoras escaped persecution in 510 B. C, though his followers were not so fortunate. D'Olivet incurred the enmity of Napoleon Bonaparte and was slated for deportation to Africa. Deportation still being in vogue, Miss Redfield may suffer a like fate, especially if the United States Senate ever delves into the Hermeneutic Interpretation of the Origin of the Social State of Man, wherein D'Olivet declared that this government is not merely indifferent to religion, but actually atheistic. No Senator who has stood up for two, four or six years under the invocations of the chaplain and his guests of all religious faiths will tolerate that accusation.

D'Olivet is not amiably Inclined toward republics. Although the United States at the time of his writings was barely in existence, he prophesied that liberty—which he considers a vital force in spiritual and political evolution—would not long endure here. He warns Europeans, inspired to hopes of a republic by America's example, that ``such a republic cannot belong to Europe unless Europe consents to become the conquest of America and to be one of its dependencies." And if America becomes strong enough to attempt such conquests he prophesies its overthrow.

The reason for D'Olivet's gloomy outlook on republicanism lies in his examination of the history of the Borean, or white, race for 12,000 years. The conclusions he was forced to draw from that monumental study agree with his interpretation of the Pythagorean cult, of which he was an Initiate.

In brief, ho perceived that there was a metaphysical correlation of Providence, Destiny and the Will of Man, which is working toward the complete evolution of man until, ``ascending into radiant ether, midst the Immortals, thou shalt be thyself a God." Destiny, according to D'Olivet, reigns over the past, the Will of Man over the future and Providence over the present.

If, he reasons, man belongs to Destiny, ``he would be what short-sighted philosophers have attributed him to be, without progression in his course and consequently without future." But, as the work of Providence always mitigates destiny, man advances freely in the route which is traced for him, perfecting himself in proportion as he advances and tending thus to immortality. Providence he conceives as the expression of the Divine Will, and Destiny the domain of the individual, in which domain, however, man has the power to dominate and regulate conditions according to the efficiency of his will.

In tracing the history of mankind D'Olivet follows the course of the ancient religions and philosophies which, to him, are the mainsprings of progress toward the ultimate state of man as a god. Without consideration of the sacred books of the nation, he declares that history is meaningless.

However one may question his conclusions, the study of history in the light of the three principles which govern his philosophy is intensely interesting. Even fine print, which is the curse of the reviewer and sudden death to the casual searchings of the layman into philosophy, could not hide the fascination of these two books. But the fact that they are beautifully bound and printed in type that can be read as easily as the best printed novel is cause not merely for honorable mention but whole-hearted rejoicing. The English translation, too, should come in for its share of commendation because of its clarity and simplicity, an end not lightly to be achieved in works of this character.

D'Olivet's interpretation of the Golden Verses, not being controversial—or perhaps we should say less so, for here we tread on dangerous ground—is more in line with popular aspirations of the moment. We may differ on politics, but, if no other sign were pointing in that direction, the number of books on the subject which pour out each month would convince us that the whole world is seeking for spiritual light. It may be rash to say we are in for a spiritual regeneration but at least it is safe to say that we are willing to be regenerated. If we don't need it ourselves, the other fellow does.

It is strange to think that seventy-one lines of verse, which are all that constitute the Golden Verses of Pythagoras, should have such tremendous effect as these esoteric verses have had in the world since the time they were written, more than 500 years before Christ. The secret of tolerance which the followers of Pythagoras, in common with the followers of Kong-Tse, or Confucius, have shown toward the gods of others is here revealed. Behind polytheism, or, indeed, the worship of any god or gods, was their belief in the essential unity of God. Consequently it didn't matter which quality the various cults deified. Pythagoreans could publicly worship the gods of the country in which they lived and know that they were really worshiping the Supreme Deity,

The laws that Pythagoras laid down for his followers were indicative of the highest standards. Since they were taught in secret, and their significance revealed only after years of probation, his doctrines were not for the general public, which is a point against the Pythagoreans and one for the Christians. Essentially, however, the moral teachings of Christ do not differ.

Pythagoras used the language of numbers (which was well known to the ancients but perfected by him) to conceal his teachings. D'Olivet does not go into this beyond a few statements which do not shed much light on the symbols.

Our chief criticism of this volume is that so much of it (the first half) should have been devoted to a dissertation on poetry when there was so much left unsaid about Pythagoras. Not that we underestimate D'Olivet's discussion of poetry. It is a most exhaustive study and should be an eye-opener to the free versifiers of today who flatter themselves that they have discovered a new thing. But, entertaining as it is, it seems to us that it has no place in the volume. Beyond dividing the book in two we fear that the publishers have no choice but to print it as it is—that is, unless D'Olivet, who was a firm believer in the transmigration of souls in their progress toward immortality, should return to earth and reveal many things concerning Pythagoras which he has concealed.

There is one point in particular on which we would like to be enlightened. Pythagoras outlined a diet for his followers which D'Olivet mentions, hut does not give in detail. Beans, however, were a Pythagorean aversion. We have a similar aversion. We recall the time at the distant age of seven when we went on a hunger strike for twenty four hours—the first hunger strike on record, wo believe. Our parents were Christians, but they were Pythagorean in that they respected other folks' principles. We can't help regretting that we didn't have a principle—an esoteric one at that—for them to respect in the matter of beans.

\flrightit{Posted on 2011-06-21 by Cologero }
